% W4i30_P_updates.tex
% DFD 17-Agent Research Programme — Iteration 30 (i30)
% Agents 6 (P1-P4), 7 (P5-P7/ET/CE), 8 (P8-P11/SQMS)
% Date: 2026-03-21
% ITERATION 18 of 20 — SYNTHESIS PHASE

\documentclass[11pt,a4paper]{article}
\usepackage[utf8]{inputenc}
\usepackage{amsmath,amssymb,amsfonts}
\usepackage{hyperref}
\usepackage{booktabs}
\usepackage{longtable}
\usepackage{geometry}
\geometry{margin=2.5cm}

\title{DFD i30 Patent-Linked Updates\\Agents 6, 7, 8\\[6pt]\large Iteration 18/20 --- Synthesis Phase}
\author{DFD 17-Agent Research Programme}
\date{2026-03-21}

\begin{document}
\maketitle
\tableofcontents
\newpage

%=============================================================================
\section{Agent 6: P1--P4 --- Final Dead Patent Confirmation}
%=============================================================================

\subsection{Mandate}
Final check on P1--P4 (LPI protocol, local tests, MEMS, interferometric). Confirm dead status. Recommend cessation of monitoring.

\subsection{New Literature (i30)}

\begin{enumerate}
\item \textbf{``A Quality Framework for Testing Gravity with Wide Binaries: No Evidence for MOND''} --- Cookson \& Banik (2026). arXiv: \href{https://arxiv.org/abs/2602.24035}{2602.24035}. MNRAS 547 (2026).\\
Rigorous framework for wide binary gravity tests. Gaia DR3 sample within 130 pc, 1--30 kAU separations. No evidence for MOND-like signal after accounting for triple contamination. Relevant to P3 (DFD local test): wide binaries are insensitive to DFD's refractive field at solar-system separations. \textbf{Grade: B.}

\item \textbf{``Wide Binaries from Gaia DR3: testing GR vs MOND with realistic triple modelling''} --- Open J.\ Astrophys.\ (2025). \href{https://astro.theoj.org/article/142887}{DOI link}.\\
Updated selection cuts, FLAMES mass estimates. Triple contamination is dominant systematic. Current sensitivity insufficient to test DFD, which predicts deviations only below $a_0$ in non-screened environments. \textbf{Grade: B.}

\item \textbf{``Testing gravity with wide binaries: 3D velocities and distances''} --- A\&A (2025). \href{https://www.aanda.org/articles/aa/full_html/2025/07/aa55115-25/aa55115-25.html}{DOI link}.\\
3D velocity reconstruction from Gaia+HARPS. Reduces projection degeneracy. Still insufficient for DFD-specific tests at solar neighbourhood densities. \textbf{Grade: C.}
\end{enumerate}

\subsection{P1--P4 Final Assessment}

\begin{center}
\begin{tabular}{clll}
\toprule
\textbf{Patent} & \textbf{Description} & \textbf{Status} & \textbf{Reason} \\
\midrule
P1 & LPI protocol & DEAD (RED) & Superseded by scalar-refractive terminology \\
P2 & Local equivalence test & DEAD (RED) & MICROSCOPE $10^{-15}$ precludes detection \\
P3 & MEMS accelerometer & DEAD (RED) & Below detection floor \\
P4 & Interferometric test & DEAD (RED) & No viable experimental path \\
\bottomrule
\end{tabular}
\end{center}

\textbf{RECOMMENDATION:} Cease monitoring P1--P4 after i30. Wide binary tests confirm that local gravity tests at these scales are insensitive to DFD. DFD's screening mechanism ($n = e^\psi$ with $\psi \ll 1$ in the solar system) ensures deviations are below $10^{-15}$, consistent with MICROSCOPE's null result.

\textbf{No further literature monitoring required.}

%=============================================================================
\section{Agent 7: P5--P7 --- ET, CE, and Next-Generation GW Detectors}
%=============================================================================

\subsection{Mandate}
Einstein Telescope site decision update. Cosmic Explorer status. DFD GW predictions for 3G detectors.

\subsection{New Literature (i30)}

\begin{enumerate}
\item \textbf{Einstein Telescope: Sardinia--Saxony Declaration of Intent (January 2026).} \href{https://www.einstein-telescope.it/en/2026/01/12/einstein-telescope-a-declaration-of-intent-between-sardinia-and-saxony-to-strengthen-scientific-cooperation-signed/}{ET Italia link}.\\
Sardinia and Saxony signed cooperation agreement for ``double-L'' configuration: two L-shaped interferometers ($\sim$15 km arms each) at two geographically distant sites. This is a \textbf{major design shift} from the original single-triangle concept. \textbf{Grade: A.}

\item \textbf{Einstein Telescope: Sardinia--Saxony dialogue continues (March 2026).} \href{https://www.einstein-telescope.it/en/2026/03/06/einstein-telescope-the-dialogue-between-italy-and-saxony-continues/}{ET Italia link}.\\
Ongoing negotiations. Double-L configuration endorsed as most cost-effective and scientifically productive. Netherlands/Belgium/Germany Euregio Meuse-Rhine remains third candidate site. \textbf{Grade: B.}

\item \textbf{Cosmic Explorer: Site evaluation and cost estimation.} arXiv: \href{https://arxiv.org/abs/2509.05430}{2509.05430} (2025).\\
26 candidate US locations for 40 km CE arms. NSF awarded \$9M for design. Preliminary site report due Fall 2026; short list 2028. Operations mid-2030s to 2040s. \textbf{Grade: B.}

\item \textbf{LVK: IR1 observing run planned for 2026--2027.} \href{https://observing.docs.ligo.org/plan/}{LVK plan}.\\
Six-month run beginning late summer/early fall 2026. A\# upgrades may be incorporated. Bridge between O4 and O5. \textbf{Grade: B.}

\item \textbf{``Forecasting Constraints on Modified GW Propagation by Strongly Lensed GWs''} --- (2026). arXiv: \href{https://arxiv.org/abs/2601.21820}{2601.21820}.\\
Forecasts for testing modified GW propagation with ET+CE using lensed events. Sensitive to $c_T \neq c$ at $\delta c_T / c \sim 10^{-15}$. DFD predicts $c_T = c$ (tensor speed unmodified), so this is a null test for DFD but a powerful discriminator against alternatives. \textbf{Grade: A.}
\end{enumerate}

\subsection{ET Timeline Update}

\begin{center}
\begin{tabular}{ll}
\toprule
\textbf{Milestone} & \textbf{Date} \\
\midrule
ET-ERIC established & 2024 \\
Sardinia--Saxony cooperation & January 2026 \\
Bid book (all 3 sites) & December 2026 \\
Site decision & H2 2027 \\
Construction begins & 2028--2029 \\
First operations & 2035 \\
\bottomrule
\end{tabular}
\end{center}

\subsection{DFD GW Predictions for 3G Detectors}

\begin{center}
\begin{tabular}{lll}
\toprule
\textbf{Prediction} & \textbf{DFD Value} & \textbf{3G Sensitivity} \\
\midrule
$c_T / c$ & $= 1$ (exact) & $\delta c_T / c \sim 10^{-15}$ (ET+CE) \\
GW luminosity distance & Modified by $\psi$ & $\delta d_L / d_L \sim 0.1\%$ (ET) \\
GW lensing magnification & $n$-dependent & Detectable with $\sim 100$ lensed events \\
Scalar GW mode (breathing) & $\lesssim 10^{-3}$ of tensor & Below 3G sensitivity \\
\bottomrule
\end{tabular}
\end{center}

\textbf{P5--P7 Status: GREEN.} Double-L ET configuration actually \emph{improves} DFD testability (better sky localisation for lensing). CE site selection advancing. DFD's $c_T = c$ prediction is a powerful null test.

%=============================================================================
\section{Agent 8: P8--P11 --- SQMS White Paper Draft}
%=============================================================================

\subsection{Mandate}
SQMS 2.0 white paper draft for internal review committee. 2--3 page document. Parasitic search mode for DFD scalar field coupling.

\subsection{New Literature (i30)}

\begin{enumerate}
\item \textbf{``Search for Postinflationary QCD Axions with a Quantum-Limited Tunable Microwave Receiver''} --- Sardo Infirri et al.\ Phys.\ Rev.\ Lett.\ 135, 211002 (2025).\\
SQMS Dark SRF first physics result. World-best constraint on dark photon kinetic mixing at specific mass range. Quality factor $Q \sim 10^{10}$. Demonstrates the sensitivity level relevant for DFD scalar coupling search. \textbf{Grade: A.}

\item \textbf{``Discovery of Niobium Hydride Precipitates in Superconducting Qubits''} --- Sung et al.\ Phys.\ Rev.\ Materials 10, 016201 (2026).\\
Materials science advance. NbH precipitates identified as decoherence source. Mitigation improves $Q$ factors further, directly benefiting Dark SRF sensitivity. \textbf{Grade: B.}

\item \textbf{``Fermilab-led study advances development of sophisticated quantum sensors''} --- Fermilab News (March 2026). \href{https://news.fnal.gov/2026/03/fermilab-led-study-advances-development-of-sophisticated-quantum-sensors-to-track-high-energy-particles-and-detect-dark-matter/}{Link}.\\
Superconducting microwire single photon detectors (SMSPDs). Complementary technology for DFD scalar field detection. \textbf{Grade: B.}

\item \textbf{``Fermilab New Technologies for Axion and Dark Photon Searches''} --- Fermilab-PUB-24-0933-T (2024). \href{https://lss.fnal.gov/archive/2024/pub/fermilab-pub-24-0933-t.pdf}{PDF}.\\
Comprehensive review of SRF cavity technologies for BSM searches. Design sensitivity $g_{a\gamma\gamma} \sim 2 \times 10^{-11}$ GeV$^{-1}$ expected by 2025--2026. Current achieved: $g_{a\gamma\gamma} \lesssim 6 \times 10^{-10}$ GeV$^{-1}$. \textbf{Grade: A.}
\end{enumerate}

\subsection{SQMS White Paper --- Draft Elements}

The following constitutes the core content for a 2--3 page SQMS Internal Review Committee white paper:

\subsubsection{Title}
``Parasitic Search for Scalar-Refractive Field Coupling Using SQMS Dark SRF Infrastructure''

\subsubsection{1. Executive Summary (0.5 page)}

Density Field Dynamics (DFD) predicts a scalar refractive field $\psi$ with $n = e^\psi$ that modifies the local speed of light as $c_\mathrm{local} = c \cdot e^{-\psi}$. This field couples to the electromagnetic sector via the refractive index, producing a frequency-dependent shift in cavity resonance frequencies. The SQMS Dark SRF infrastructure---with $Q \sim 10^{10}$ niobium cavities at 2K---is uniquely positioned to detect or constrain this coupling in a parasitic search mode that does not interfere with the primary dark photon/axion programme.

\subsubsection{2. Physics Case (0.5 page)}

\begin{itemize}
\item DFD coupling: $\delta\nu / \nu \sim \psi \sim G M / (r c^2) \sim 10^{-9}$ (Earth surface).
\item Time variation: $\dot{\psi} / \psi \sim H_0 \sim 10^{-18}$ s$^{-1}$ (cosmological).
\item Modulation: Annual variation from Earth's orbital eccentricity: $\delta\psi_\mathrm{annual} \sim 3 \times 10^{-10}$.
\item Target sensitivity: $\delta\nu / \nu \sim 10^{-18}$ per $\sqrt{\text{Hz}}$ (achievable with $Q \sim 10^{10}$).
\item Discriminator: DFD predicts \emph{frequency-independent} fractional shift (unlike axion, which is mass-dependent).
\end{itemize}

\subsubsection{3. Experimental Approach (0.5 page)}

\begin{itemize}
\item \textbf{Primary observable:} Long-term ($\sim 1$ year) monitoring of cavity resonance frequency with $10^{-18}$ fractional stability.
\item \textbf{Modulation search:} Annual sinusoidal modulation at orbital period (365.25 days) with amplitude $\delta\nu / \nu \sim 3 \times 10^{-19}$.
\item \textbf{Parasitic mode:} Use existing Dark SRF data. No additional beam time required. Analysis-only cost.
\item \textbf{Systematics:} Temperature stability (already controlled at mK level). Pressure variations (cryostat isolated). Magnetic field (shielded).
\item \textbf{Cross-check:} Simultaneous monitoring of two cavities at different frequencies provides frequency-independence test.
\end{itemize}

\subsubsection{4. Resource Requirements (0.25 page)}

\begin{itemize}
\item \textbf{Hardware:} None beyond existing Dark SRF infrastructure.
\item \textbf{Personnel:} 0.5 FTE postdoc for 12 months (data analysis).
\item \textbf{Computing:} Modest ($\sim 10^4$ CPU-hours for time-series analysis).
\item \textbf{Total cost:} \$75k--\$100k (postdoc salary + overhead).
\end{itemize}

\subsubsection{5. Expected Outcome (0.25 page)}

\begin{itemize}
\item \textbf{Best case:} Detection of annual modulation consistent with DFD. First experimental evidence for scalar-refractive coupling.
\item \textbf{Null result:} Constraint on DFD scalar coupling strength at $\psi < 10^{-10}$ level, improving on current laboratory bounds by $\sim 2$ orders of magnitude.
\item \textbf{Publication:} 1 PRL-class paper regardless of outcome.
\end{itemize}

\subsection{SQMS Review Committee Submission Path}

\begin{itemize}
\item \textbf{Step 1:} Format as SQMS-NOTE document (April 2026).
\item \textbf{Step 2:} Internal circulation to Dark SRF team lead (Romanenko).
\item \textbf{Step 3:} Presentation at SQMS 2.0 kickoff meeting (expected Q3 2026).
\item \textbf{Step 4:} If endorsed, formal proposal to SQMS Internal Review Committee.
\end{itemize}

\textbf{P8--P11 Status: GREEN.} SQMS 2.0 funded. White paper draft ready for formatting. Parasitic search mode minimises resource competition. PRL 135, 211002 (2025) demonstrates the sensitivity floor.

%=============================================================================
\section{Patent Agents --- Combined i30 Assessment}
%=============================================================================

\begin{center}
\begin{tabular}{clll}
\toprule
\textbf{Agent} & \textbf{Patents} & \textbf{Status} & \textbf{Action} \\
\midrule
6 & P1--P4 & RED (dead) & Cease monitoring after i30 \\
7 & P5--P7 & GREEN & ET double-L; CE site eval; IR1 2026 \\
8 & P8--P11 & GREEN & SQMS white paper ready for formatting \\
\bottomrule
\end{tabular}
\end{center}

\end{document}
